\subsection{Poincaré series of a graded module over a polynomial ring}

Let \(H_{0}\) be a ring, I a finite set, and H the polynomial ring \(\mathrm{H}_{0}\big[(\mathrm{X}_{i})_{i\in\mathbf{I}}\big]\) . For each \(i\in I\) , let \(d_{i}\) be an integer \(>0\). Let us endow H with the structure of a graded ring of type Z such that the elements of \(H_{0}\) are homogeneous of degree 0 and each \(X_{i}\) is homogeneous of degree \(d_{i}\) . When \(d_{i}=1\) for every \(i\) , we recover the usual grading of polynomial rings.

Let M be a finitely generated graded H-module all of whose homogeneous components are \(H_{0}\) -modules of finite length; one calls the Poincaré series of M the element \(P_{M}\) of \(\mathbf{Z}((T))\) such that \(\mathrm{P}_{\mathrm{M}} = \sum_{n \in \mathbf{Z}} \mathrm{long}_{\mathrm{H}_{0}}(\mathrm{M}_{n}) \cdot \mathrm{T}^{n}\) , and one sets \(Q_{M} = P_{M} \cdot \prod_{i \in I} (1 - T^{d_{i}})\) .

THEOREM 1. --- The element \(\mathbf{Q}_{\mathbf{M}}\) of \(\mathbf{Z}((\mathbf{T}))\) belongs to \(\mathbf{Z}[\mathbf{T}, \mathbf{T}^{-1}]\) .

Dividing \(H_{0}\) by the annihilator of the \(H_{0}\) -module M, we reduce to the case where M is a faithful \(H_{0}\) -module. If \(a, b \in Z\) are such that M is generated as an H-module by \(M' = \sum_{a \leqslant i \leqslant b} M_{i}\) , then \(M'\) is a nonzero finitely generated \(H_{0}\) -module faithful and of finite length; consequently,

the ring \(\mathbf{H}_0\) is Artinian (A, VIII, § 1, no. 3), hence Noetherian (loc. cit., § 9, no. 1). The polynomial ring H is therefore Noetherian (loc. cit., § 1, no. 4). If I is empty, we have \(\mathbf{H} = \mathbf{H}_0\) , and the family of integers \((\mathrm{long}_{\mathbf{H}_0}(\mathbf{M}_n))_{n \in \mathbf{Z}}\) has finite support, since M is a finitely generated \(\mathbf{H}_0\) -module; whence the theorem in this case. We then argue by induction on the cardinality of the set I, assumed nonempty; let \(j \in \mathbf{I}\) and \(\mathbf{J} = \mathbf{I} - \{j\}\) . Denote by H' the graded subring of H generated by \(\mathbf{H}_0\) and the \(X_i\) for \(i\) in J; consider the homothety \((\mathbf{X}_j)_\mathbf{M}\) of ratio \(X_j\) on M, its kernel R, and its cokernel S. For each \(n \in \mathbf{Z}\) , we have an exact sequence of \(\mathbf{H}_0\) -modules

\[
0 \rightarrow \mathrm{R} _ {n - d _ {j}} \rightarrow \mathrm{M} _ {n - d _ {j}} \rightarrow \mathrm{M} _ {n} \rightarrow \mathrm{S} _ {n} \rightarrow 0,
\]

hence \(\mathbf{R}_{n - d_j}\) and \(\mathbf{S}_n\) are of finite length, and we have

\[
\operatorname{long} _ {\mathrm{H} _ {0}} \left(\mathrm{M} _ {n}\right) - \operatorname{long} _ {\mathrm{H} _ {0}} \left(\mathrm{M} _ {n - d _ {j}}\right) = \operatorname{long} _ {\mathrm{H} _ {0}} \left(\mathrm{S} _ {n}\right) - \operatorname{long} _ {\mathrm{H} _ {0}} \left(\mathrm{R} _ {n - d _ {j}}\right).\tag{4}
\]

Since M is a finitely generated module over the Noetherian ring H, the H-modules R and S are finitely generated; as they are annihilated by \(X_{j}\) , they are finitely generated H'-modules. By the induction hypothesis, the elements \(P_{R}\cdot\prod_{i\in J}(1-T^{d_{i}})\) and \(P_{S}\cdot\prod_{i\in J}(1-T^{d_{i}})\) of \(\mathbf{Z}((T))\) therefore belong to \(\mathbf{Z}[T,T^{-1}]\) ; from (4), we deduce

\[
\mathrm{P} _ {\mathrm{M}} - \mathrm{T} ^ {d _ {j}}. \mathrm{P} _ {\mathrm{M}} = \mathrm{P} _ {\mathrm{S}} - \mathrm{T} ^ {d _ {j}}. \mathrm{P} _ {\mathrm{R}},
\]

that is, \((1 - \mathrm{T}^{d_j}).\mathrm{P}_{\mathrm{M}} = \mathrm{P}_{\mathrm{S}} - \mathrm{T}^{d_j}.\mathrm{P}_{\mathrm{R}}\) ; we thus have

\[
\mathrm{P} _ {\mathrm{M}} \cdot \prod_ {i \in \mathrm{I}} (1 - \mathrm{T} ^ {d _ {i}}) = \mathrm{P} _ {\mathrm{S}} \cdot \prod_ {i \in \mathrm{J}} (1 - \mathrm{T} ^ {d _ {i}}) - \mathrm{T} ^ {d _ {j}} \cdot \mathrm{P} _ {\mathrm{R}} \cdot \prod_ {i \in \mathrm{J}} (1 - \mathrm{T} ^ {d _ {i}}),\tag{5}
\]

whence the conclusion.

Example 1. --- Suppose \(H_{0}\) is Artinian and take M = H. Then, with the previous notation, we have R = 0 and S = H', so by (5), we have \(Q_{H} = Q_{H'}\) ; since we have \(Q_{H_{0}} = \text{long}(H_{0})\) from which we obtain by induction \(Q_{H} = \text{long}(H_{0})\) , that is,

\[
\mathrm{P} _ {\mathrm{H}} = \operatorname{long} (\mathrm{H} _ {0}). \prod_ {i \in \mathrm{I}} (1 - \mathrm{T} ^ {d _ {i}}) ^ {- 1}.\tag{6}
\]

Suppose henceforth that H is equipped with the usual grading, for which \(d_{i}=1\) for every \(i\in I\) , and set \(r=\operatorname{Card}(\mathbf{I})\) ; one has \(\mathrm{P}_{\mathrm{M}}=\mathrm{Q}_{\mathrm{M}}(\mathrm{T}).(1-\mathrm{T})^{-r}\) . Put \(c_{\mathrm{M}}=\mathrm{Q}_{\mathrm{M}}(1)\) . Then, by formula (2) of no. 1, we have:

COROLLARY. --- a) If r = 0, then we have long \(_{H_{0}}(M)\) = c \(_{M}\) .

\begin{enumerate}
\def\labelenumi{\alph{enumi})}
\setcounter{enumi}{1}
\item
  If \(r = 1\) , then we have \(\mathrm{long}_{\mathrm{H}_0}(\mathbf{M}_n) = c_{\mathbf{M}}\) for \(n\) sufficiently large.
\item
  If \(r > 1\) , then we have long \(_{\mathrm{H}_0}(\mathbf{M}_n) = c_{\mathbf{M}} \frac{n^{r-1}}{(r - 1)!} + \rho_n n^{r-2}\) , where \(\rho_n\) tends to a limit in \(\mathbf{R}\) when \(n\) increases indefinitely.
\end{enumerate}

Remarks. --- 1) The integer \(c_{\mathbf{M}}\) is positive by Lemma 2. One can have \(\mathbf{M} \neq 0\) and \(c_{\mathbf{M}} = 0\) (cf. prop. 2).

\begin{enumerate}
\def\labelenumi{\arabic{enumi})}
\setcounter{enumi}{1}
\tightlist
\item
  Let \(0 \to M' \to M \to M'' \to 0\) be an exact sequence of graded H-modules and degree-0 homomorphisms such that M is finitely generated over H and \(M_n\) is of finite length over \(H_0\) for each \(n\). Then, for each \(n \in \mathbf{Z}\) , one has
\end{enumerate}

\[
\operatorname{long} _ {\mathrm{H} _ {0}} \left(\mathrm{M} _ {n}\right) = \operatorname{long} _ {\mathrm{H} _ {0}} \left(\mathrm{M} _ {n} ^ {\prime}\right) + \operatorname{long} _ {\mathrm{H} _ {0}} \left(\mathrm{M} _ {n} ^ {\prime \prime}\right),
\]

hence \(\mathrm{P_M} = \mathrm{P_{M'}} + \mathrm{P_{M''}}\) , \(\mathrm{Q_M} = \mathrm{Q_{M'}} + \mathrm{Q_{M''}}\) and \(c_{\mathbf{M}} = c_{\mathbf{M'}} + c_{\mathbf{M''}}\) .

\begin{enumerate}
\def\labelenumi{\arabic{enumi})}
\setcounter{enumi}{2}
\tightlist
\item
  Let \(\mathbf{M}(p)\) be the module deduced from \(\mathbf{M}\) by shift by \(p\) in the grading (A, II, p. 165, example 3). Since we have \(\mathbf{M}(p)_n = \mathbf{M}_{p+n}\) , one has \(\mathbf{P}_{\mathbf{M}(p)} = \mathrm{T}^{-p}\mathbf{P}_{\mathbf{M}}\) , \(\mathbf{Q}_{\mathbf{M}(p)} = \mathrm{T}^{-p}\mathbf{Q}_{\mathbf{M}}\) and \(c_{\mathbf{M}(p)} = c_{\mathbf{M}}\) .
\end{enumerate}

Examples. --- 2) Suppose H \(_{0}\) is Artinian, and let M be a graded free H-module generated by s homogeneous, linearly independent elements of respective degrees \(\delta_{1}, ..., \delta_{s}\). Then M is isomorphic to H( \(-\delta_{1}\) ) ⊕ \(\cdots\) ⊕ H( \(-\delta_{s}\) ). By Remarks 2 and 3 and Example 1, we therefore have

\[
\begin{array}{l} \mathrm {P_ {M}} = \operatorname{long} (\mathrm {H_ {0}}) \left(\sum_ {i = 1} ^ {s} \mathrm{T} ^ {\delta_ {i}}\right) (1 - \mathrm{T}) ^ {- r}, \\ \mathrm {Q_ {M}} = \operatorname{long} (\mathrm {H_ {0}}) \left(\sum_ {i = 1} ^ {s} \mathrm{T} ^ {\delta_ {i}}\right), \\ c _ {\mathrm{M}} = s. \operatorname{long} (\mathrm {H_ {0}}). \end{array}
\]

\begin{enumerate}
\def\labelenumi{\arabic{enumi})}
\setcounter{enumi}{2}
\tightlist
\item
  Suppose again that \(H_{0}\) is Artinian; let M be a graded H-module, and suppose there exists an exact sequence of graded H-modules and degree-0 homomorphisms
\end{enumerate}

\[
0 \to \mathrm{L} _ {n} \to \mathrm{L} _ {n - 1} \to \dots \to \mathrm{L} _ {0} \to \mathrm{M} \to 0,
\]

where, for \(k=0,1,...,n\) , \(\mathbf{L}_{k}\) is a free graded H-module generated by linearly independent homogeneous elements, of respective degrees \(\delta_{k,1},..., \delta_{k,m(k)}\). Then, by Remark 2 and Example 2, we have

\[
\mathrm{Q} _ {\mathrm{M}} = \operatorname{long} (\mathrm{H} _ {0}). \sum_ {0 \leqslant k \leqslant n} \sum_ {1 \leqslant j \leqslant m (k)} (- 1) ^ {k} \mathrm{T} ^ {\delta_ {k, j}},
\]

\[
c _ {\mathbf {M}} = \operatorname{long} (\mathrm{H} _ {0}). \sum_ {0 \leqslant k \leqslant n} (- 1) ^ {k} m (k).
\]

Remark 4. --- One can prove (p. 88, exerc. 4) that under the hypotheses of th. 1, the \(\mathrm{H}_{0}\) -modules \(\mathrm{Tor}_{j}^{\mathrm{H}}(\mathrm{H}_{0},\mathbf{M})\) are of finite length, are zero for \(j > r\) , and one has \(c_{\mathbf{M}} = \sum_{j=0}^{r} (-1)^{j} \mathrm{long}_{\mathrm{H}_{0}}(\mathrm{Tor}_{j}^{\mathrm{H}}(\mathrm{H}_{0},\mathbf{M}))\). More precisely, the H-modules

\(T_{j} = \text{Tor}_{j}^{\text{H}}(\text{H}_{0}, \text{M})\) are naturally equipped with graduations, and one has

\[
\mathrm{Q} _ {\mathrm{M}} = \sum_ {j = 0} ^ {r} (- 1) ^ {j} \mathrm{P} _ {\mathrm{T} _ {j}}.
\]

PROPOSITION 1.--- Let M be a graded H-module. Suppose that M is generated by \(M_{0}\) and that \(M_{0}\) is an \(H_{0}\)-module of finite length. Then we have

\[
\mathrm{P} _ {\mathrm{M}} \leqslant (1 - \mathrm{T}) ^ {- r} \operatorname{long} _ {\mathrm{H} _ {0}} (\mathrm{M} _ {0}), \quad c _ {\mathrm{M}} \leqslant \operatorname{long} _ {\mathrm{H} _ {0}} (\mathrm{M} _ {0}).
\]

Moreover, the following conditions are equivalent:

\begin{enumerate}
\def\labelenumi{(\roman{enumi})}
\tightlist
\item
  \(c_{\mathbf{M}} = \mathrm{long}_{\mathrm{H_0}}(\mathbf{M}_0)\) ;
\end{enumerate}

\[
\text{(ii)} \mathrm{P}_{\mathbf{M}} = \operatorname{long}_{\mathbf{H}_{0}}(\mathbf{M}_{0}) \cdot (1 - \mathrm{T})^{-r},\ \text{that is, }\mathbf{M} = \mathbf{M}_{0}\text{ if }r=0\text{ and}
\]

\[
\operatorname{long} _ {\mathrm{H} _ {0}} \left(\mathrm{M} _ {n}\right) = \operatorname{long} _ {\mathrm{H} _ {0}} \left(\mathrm{M} _ {0}\right). \binom {n + r - 1} {r - 1}
\]

for \(n \in \mathbf{N}\) if \(r > 0\) ;

\begin{enumerate}
\def\labelenumi{(\roman{enumi})}
\setcounter{enumi}{2}
\tightlist
\item
  the canonical homomorphism of H-modules
\end{enumerate}

\[
\varphi : \mathrm{H} \otimes_ {\mathrm{H} _ {0}} \mathrm{M} _ {0} \rightarrow \mathrm{M}
\]

is bijective.

Let R denote the kernel of φ. Since φ is surjective, we have

\(P_{M} = P_{H \otimes M_{0}} - P_{R} = \text{long}_{H_{0}}(M_{0})(1 - T)^{-r} - P_{R}\) and \(c_{M} = \text{long}_{H_{0}}(M_{0}) - c_{R}\). The conditions (i), (ii), and (iii) are respectively equivalent to \(c_{R} = 0\), \(P_{R} = 0\) and \(R = 0\). We therefore have (iii) \(\Rightarrow\) (ii) \(\Rightarrow\) (i), and it suffices to prove that \(c_{R} = 0\) implies \(R = 0\). Suppose \(R \neq 0\) and let \(0 = N^{h} \subset N^{h-1} \subset \ldots \subset N^{0} = M_{0}\) be a Jordan–Hölder sequence of the \(H_{0}\)-module \(M_{0}\). Let \(R^{m}\) be the intersection of \(R\) and the image of \(H \otimes_{H_{0}} N^{m}\) in \(H \otimes_{H_{0}} M_{0}\); there exists an integer \(m\) between 1 and \(h\) such that \(R^{m} \neq R^{m-1}\). Put \(L = R^{m-1}/R^{m}\); one has \(0 \leqslant c_{L} \leqslant c_{R}\) and it suffices to prove that \(c_{L} \neq 0\). Now, if \(k\) is the quotient field of \(H_{0}\) by the maximal annihilator ideal of \(N^{m-1}/N^{m}\), L is identified with a nonzero graded submodule of \(k[(X_{i})_{i\in I}]\). Therefore L contains a submodule isomorphic to a shifted module of \(k[(X_{i})_{i\in I}]\); since \(c_{k[(X_{i})_{i\in I}]} = 1\), one therefore has \(c_{L} \geqslant 1\) (remarks 2 and 3), which is what we wanted to prove.

By A, X, p. 160, th. 1, condition (iii) means that \((\mathbf{X}_{1}, \ldots, \mathbf{X}_{r})\) is a completely secant sequence for the H-module M.

PROPOSITION 2. --- Suppose that \(H_{0}\) is a field, and let M be a finitely generated graded H-module. Let K be the field of fractions of H. Then \(c_{M}\) is equal to the rank of the H-module M, that is, to the dimension of the K-vector space \(M \otimes_{H} K\) .

This is clear if M = H, since \(c_{H} = 1\) . Moreover, let \(x \in H\) , homogeneous of degree d, and nonzero; we have (H/xH) \(\otimes_{H} K = 0\) ; from the exact sequence

\[
0 \rightarrow \mathrm{H} (- d) \rightarrow \mathrm{H} \rightarrow \mathrm{H} / x \mathrm{H} \rightarrow 0,
\]

and Remarks 2 and 3, we deduce \(c_{\mathrm{H}/x\mathrm{H}}=0\) . The proposition is therefore verified when M is generated by a single homogeneous element. The general case follows from this, since every finitely generated graded H-module has a composition series whose quotients are of the preceding form.

Remark 6. -- Under the hypotheses of Prop. 2, we therefore have \(c_{\mathbf{M}} = 0\) if and only if M is a torsion H-module, or equivalently if and only if \(\dim_{\mathrm{H}}(\mathbf{M}) < r\) ( \(\S 1, \mathfrak{n}^{\circ} 5\) , Example 4).

